% Source audit, JAC1964-20261121; /root/programme_audit.
% This note does not replace or strengthen JAC-P16 unconditionally.
\subsection*{JAC-P20. An actual abelian cohomology injection and the separate transfer audit}

Milne, \emph{Grothendieck's standard conjecture of Lefschetz type
over finite fields}, version 1.1 of 14 July 2022, Proposition~4.2,
printed pp.~7--8, explicitly asserts the following transfer:
if an algebraic correspondence induces an injection
$H^*(X)\longrightarrow H^*(Y)$ and the Hodge standard conjecture
holds for $Y$, then it holds for $X$.
Its proof consists of references to Kleiman (1968), Theorem~3.11,
and Saavedra Rivano (1972), VI, 4.4.2.
The proposition is a stated result in the primary source.
The present audit has \emph{not re-established} its positivity
transfer from those two cited passages. This is not a refutation
of the proposition.

\begin{lemma}[JAC-P20]\label{jac:P20}
Let $X$ be a smooth projective variety over an algebraically
closed field whose rational Chow motive is of abelian type.
There are an abelian variety $Y$ and an actual rational algebraic,
possibly inhomogeneous correspondence on $X\times Y$ that induces
an injection $H^*(X)\longrightarrow H^*(Y)$ in every classical
Weil cohomology theory. Here total cohomology is regarded without
its degree and Tate-twist labels.
\end{lemma}

\begin{proof}
Choose an actual Chow-motive splitting
\[
 h(X)\xrightarrow{u}\bigoplus_{i=1}^{s}h(A_i)(n_i)
       \xrightarrow{v}h(X),\qquad vu=1,
\]
where the $A_i$ are abelian varieties. Such a finite presentation
exists:
products are abelian varieties, duals introduce only twists,
and Artin motives over the algebraic closure are sums of units.
Write $u_i$ for the components of $u$. Forgetting the degree
labels and Tate factors, their algebraic correspondences give
an injection into $\bigoplus_iH^*(A_i)$. The splitting identity
proves injectivity; no realization-only isomorphism is substituted
for a correspondence.

Choose an elliptic curve $E$ and set
\[
                 Y=\left(\prod_{i=1}^{s}A_i\right)\times E^s.
\]
For each $i$, define an algebraic map on total cohomology
\[
 j_i(a)=\operatorname{pr}_{A_i}^{*}a\,
                   \operatorname{pr}_{E_i}^{*}[0].
\]
All unspecified factors carry their unit class. Pullback and
multiplication by the point class are algebraic correspondences.
Each $j_i$ is injective. The images for different $i$ lie in
distinct Kunneth multidegrees of the $E^s$ factor: degree two in
the $i$th elliptic factor and degree zero in every other elliptic
factor. They therefore have zero mutual intersection. Consequently
\[
                    \alpha_* =\sum_i j_i(u_i)_*
                    :H^*(X)\longrightarrow H^*(Y)
\]
is injective and is induced by an actual algebraic correspondence.
The correspondence may have components in several codimensions;
this is precisely the ungraded $H^*$ notation of Milne's stated
Proposition~4.2. Kleiman (1968), Section~1.3, printed pp.~364--365,
likewise treats an arbitrary correspondence on total cohomology
and then separates its homogeneous components.
\end{proof}

\begin{corollary}[JAC-P20, conditional transfer]
Take Milne's Proposition~4.2, in its stated total-cohomology
form, as an additional input. Then the accepted Hodge-standard
statement of 032 for abelian varieties implies the Hodge standard
conjecture for every smooth projective variety of geometric
abelian type, in every characteristic.
\end{corollary}
\begin{proof}
Apply the additional transfer assertion to the actual injection
of JAC-P20 and to the abelian variety $Y$, for which the accepted
032 input supplies the Hodge-standard assertion.
For a nonclosed ground field the same argument is made after
algebraic closure; the intersection numbers of the original
cycles are unchanged.
\end{proof}
This corollary is conditional on that precisely identified
additional input. It is not inferred from Jannsen semisimplicity
alone and is not counted as closure of the remaining fixed-axis
sign gate.

\paragraph{What the two cited proof passages establish.}
Kleiman's Theorem~3.11, printed pp.~380--381, assumes the
Lefschetz theorem and $B(X),B(Y)$. Its first assertion says
that the transpose of an algebraic correspondence for the
\emph{geometric star forms} is algebraic and that the trace of
$u'u$ is rational. Its strict-positivity assertion additionally
assumes the Hodge standard conjecture for $X\times Y$ and the
product ample class. This is not, in the displayed statement,
positivity inferred from an injection and Hodge standard for
$Y$ alone.

Saavedra's VI, 4.2.0, printed p.~363, assumes the standard
conjectures and explicitly retains that convention in
Sections~4.3--4.6. Under that assumption, VI, 4.4.2, printed
p.~373, constructs the canonical polarization from the
geometric star forms, and verifies their mutual compatibility
using Appendix A, 3.3 and Kleiman's theorem. Appendix A, 3
starts with the same standard-conjecture assumption on printed
p.~400; its trace-positivity formula (3.3.1) is on p.~401.
These passages supply polarization formalism once positivity
of the geometric forms is available. In this audit they have
not supplied the additional argument identifying the particular
form of $X$ with a positive form transported through $\alpha_*$.

An abstract positive form on a summand does not by itself settle
that comparison: replacing it by its negative leaves the
associated adjoint involution unchanged. This observation only
identifies the missing logical step in such a proposed proof;
it is not a geometric counterexample. Accordingly JAC-P16
continues to assert its independently checked $B$ and $D$
conclusions, while the stronger conclusion above is explicitly
conditional on the additional transfer proposition.

\paragraph{Primary-source audit trail.}
Milne's original PDF was read at
\url{https://www.jmilne.org/math/articles/LFF.pdf},
Proposition~4.2 on pp.~7--8 and Theorem~4.5 on p.~8.
Kleiman's original pages were read in the Mainz scan of
\emph{Algebraic cycles and the Weil conjectures},
\emph{Dix exposes sur la cohomologie des schemas} (1968),
pp.~359--386, especially pp.~364--365 and 380--381.
Saavedra's original book, \emph{Categories tannakiennes},
Lecture Notes in Mathematics~265 (1972), was read in the
Columbia/de Jong scan at
\url{https://www.math.columbia.edu/~dejong/tannakian/},
with the exact printed locators given above. Those original
pages were inspected visually as well as by text extraction.
No erratum or contrary theorem is asserted here.
